% Level 3, Week 6 homework -- Review and Checkpoint
\homeworkweek{6}{Review and Checkpoint}{15--20 minutes. Do Core first; Challenge and Creative are optional.}
\namedateline

\begin{core}[Piggy Bank Bug Hunt (about 10 minutes)]
\begin{enumerate}[itemsep=2pt]
  \item Open \texttt{week06/hw-bug-bank/starter.py}. Read the expected output at the top of the file, then press Run. What is different?
  \item Find the 3 bugs. Use the debugger: Debug, Step over, Step into, and watch \texttt{self.coins} and \texttt{self.log} in the Variables view.
  \item Fill in your bug log, then open \texttt{check\_my\_work.py} and press Run until all checks pass. Save a copy as \texttt{bug\_bank.py} in \texttt{git-quest} (\menu{File > Save as...}), then \texttt{git add} and \texttt{git commit} it.
\end{enumerate}
\begin{center}
\renewcommand{\arraystretch}{1.6}
\begin{tabular}{|c|p{1.4cm}|p{5.2cm}|p{5.2cm}|}
\hline
\textbf{Bug} & \textbf{Line} & \textbf{What was wrong} & \textbf{How I fixed it} \\ \hline
1 & & & \\ \hline
2 & & & \\ \hline
3 & & & \\ \hline
\end{tabular}
\end{center}
\end{core}

\begin{challenge}[Trace table]
Without running it, fill in the trace table for this code. Then check with the debugger.
\begin{pycode}
stack = []
for word in ["red", "green", "blue"]:
    stack.append(word)
top = stack.pop()
\end{pycode}
\begin{center}
\renewcommand{\arraystretch}{1.4}
\begin{tabular}{|p{3cm}|p{6cm}|p{3cm}|}
\hline
\textbf{After line} & \textbf{stack} & \textbf{top} \\ \hline
append \enquote{red} & & \\ \hline
append \enquote{green} & & \\ \hline
append \enquote{blue} & & \\ \hline
\texttt{top = stack.pop()} & & \\ \hline
\end{tabular}
\end{center}
\end{challenge}

\begin{creative}[Bug comic]
Draw a 3-panel comic about a bug you met this term: panel 1, the bug appears; panel 2, the hunt (what strategy did you use?); panel 3, the fix. Bugs can be cute!
\answerspace{4cm}
\end{creative}

\begin{selfcheck}
  \item All checks pass in \texttt{hw-bug-bank}.
  \item I used the debugger at least once at home.
  \item I changed one thing at a time and ran the program after each change.
  \item I committed my fixed file in \texttt{git-quest}.
\end{selfcheck}
